%% toc: P15 166 166.4
% Proofs for relative information, divergence and sufficiency.
\section{Relative information and Kullback--Leibler divergence: proofs}
\label{sec:annex-p15-w02b-p15-045-relative-information}

Throughout, an \emph{extended integral} \(\int u\,d\rho=\int u^+\,d\rho-\int u^-\,d\rho\)
is used only when \(\int u^-\,d\rho<\infty\); its value lies in
\((-\infty,\infty]\). Two facts about such integrals are used repeatedly.

\emph{(E1) Additivity.} Let \(u,v\) be real measurable functions with
\(\int u^-\,d\rho<\infty\) and \(\int v^-\,d\rho<\infty\), and put
\(w=u+v\). Then \(w^-\leq u^-+v^-\), so \(\int w^-\,d\rho<\infty\), and
\(\int w\,d\rho=\int u\,d\rho+\int v\,d\rho\). Indeed,
\(w^+-w^-=u^+-u^-+v^+-v^-\) gives
\(w^++u^-+v^-=w^-+u^++v^+\) pointwise, all terms nonnegative.
Integrating with additivity of the nonnegative integral
(\cref{thm:p15-nonnegative-integral-laws}) gives
\(\int w^++\int u^-+\int v^-=\int w^-+\int u^++\int v^+\), and the three
finite quantities \(\int u^-,\int v^-,\int w^-\) may be subtracted.

\emph{(E2) Disjoint pieces.} If \(E_1,\dots,E_n\) are disjoint measurable
sets with union \(X\) and \(\int u^-\,d\rho<\infty\), then
\(\int u\,d\rho=\sum_j\int_{E_j}u\,d\rho\), because
\(\int u^\pm\,d\rho=\sum_j\int_{E_j}u^\pm\,d\rho\) and every
\(\int_{E_j}u^-\,d\rho\) is finite.

\emph{(P) Nonnegative multipliers.} Let \(\rho\) be a finite measure on
\(\mathbf S\), \(\mathbf S_0\subseteq\mathbf S\) a sub-\(\sigma\)-algebra,
\(V\geq0\) \(\mathbf S\)-measurable and \(V_0\geq0\)
\(\mathbf S_0\)-measurable with \(\int_AV_0\,d\rho=\int_AV\,d\rho\) for
every \(A\in\mathbf S_0\). Then \(\int WV_0\,d\rho=\int WV\,d\rho\) in
\([0,\infty]\) for every \(\mathbf S_0\)-measurable \(W\geq0\). This holds
for \(W=1_A\) by hypothesis, for nonnegative \(\mathbf S_0\)-simple \(W\)
by linearity, and for general \(W\) by taking \(\mathbf S_0\)-simple
\(0\leq W_n\uparrow W\) and applying monotone convergence
(\cref{thm:p15-monotone-integral}) to \(W_nV_0\uparrow WV_0\) and
\(W_nV\uparrow WV\). This is the nonnegative case of
\cref{thm:p15-conditional-pullout}.

\begin{annexitem}{annex:p15-w02b-p15-045-well-defined}{thm:p15-w02b-p15-045-well-defined}{The information integrals exist and do not depend on the dominating measure}
\begin{proof}
Remove from \(X\) the \(\lambda\)-null set on which \(f_1\) or \(f_2\)
lies outside \((0,\infty)\); integrals over \(\lambda\)-null sets vanish,
so nothing below changes. Put \(g=f_1/f_2\).

(1) For \(t>0\), \(t\log t\geq-1/e\), since the derivative \(1+\log t\)
vanishes only at \(t=1/e\), where the value is \(-1/e\), and
\(t\log t\to0\) as \(t\downarrow0\) and \(\to\infty\) as \(t\to\infty\).
Hence \(f_1\log(f_1/f_2)=f_2\,g\log g\geq-f_2/e\), so
\((f_1\log(f_1/f_2))^-\leq f_2/e\) and
\(\int_E(f_1\log(f_1/f_2))^-\,d\lambda\leq\mu_2(E)/e<\infty\). The
extended integral over \(E\) therefore exists in \((-\infty,\infty]\),
and dividing by \(\mu_1(E)>0\) keeps it there. Exchanging the indices
gives the same for \(I_{2:1}(E)\).

(2) For \(E\in\mathbf S\), the change-of-measure rule
(\cref{thm:p15-density-calculus}, part~1, with \(\nu=\mu_2\),
\(\mu=\lambda\), density \(f_2\) and integrand \(1_Ef_1/f_2\geq0\))
gives \(\int_E(f_1/f_2)\,d\mu_2=\int_Ef_1\,d\lambda=\mu_1(E)\). So
\(f_1/f_2\) is a version of \(d\mu_1/d\mu_2\). Let \(\lambda'\) be
\(\sigma\)-finite with \(\lambda'\equiv\{\mu_1,\mu_2\}\) and
\(f_i'=d\mu_i/d\lambda'\) (\cref{thm:p15-radon-nikodym}). Then
\(f_i'<\infty\) \([\lambda']\), and
\(\mu_i(\{f_i'=0\})=\int_{\{f_i'=0\}}f_i'\,d\lambda'=0\), so
\(\lambda'(\{f_i'=0\})=0\). The same computation shows that \(f_1'/f_2'\)
is a version of \(d\mu_1/d\mu_2\). Two versions agree \(\mu_2\)-almost
everywhere, and \(\mu_2,\lambda,\lambda'\) have the same null sets, so
\(f_1'/f_2'=f_1/f_2\) \([\lambda]\). Applying the change-of-measure rule
to the positive and negative parts of \(1_E\log(f_1/f_2)\), with
\(\nu=\mu_1\) and density \(f_1\), gives
\[
  \int_Ef_1\log\frac{f_1}{f_2}\,d\lambda=\int_E\log\frac{f_1}{f_2}\,d\mu_1,
\]
and the same identity holds for \((\lambda',f_1',f_2')\). The right-hand
sides coincide because the two ratios agree \(\mu_1\)-almost everywhere.
So \(I_{1:2}(E)\) is the same for both choices. The same argument with
the indices exchanged gives \(I_{2:1}(E)\), and hence \(J_{12}(E)\).

(3) If \(\mu_1(E)>0\) then \(\lambda(E)>0\), hence \(\mu_2(E)>0\). Put
\[
  u=1_E\frac{f_1}{\mu_1(E)}\log\frac{f_1}{f_2},\qquad
  v=-1_E\frac{f_2}{\mu_2(E)}\log\frac{f_1}{f_2}=1_E\frac{f_2}{\mu_2(E)}\log\frac{f_2}{f_1}.
\]
By (1), \(\int u^-\,d\lambda\) and \(\int v^-\,d\lambda\) are finite, and
\(\int u\,d\lambda=I_{1:2}(E)\), \(\int v\,d\lambda=I_{2:1}(E)\). By (E1),
\(\int(u+v)\,d\lambda=I_{1:2}(E)+I_{2:1}(E)=J_{12}(E)\), which is the
displayed formula. For \(E=X\), \(\mu_1(X)=\mu_2(X)=1\) and the
integrand is \((f_1-f_2)\log(f_1/f_2)\). This is nonnegative, because
\(f_1-f_2\) and \(\log(f_1/f_2)\) have the same sign. Hence
\(J(1,2)\in[0,\infty]\) and \(J(1,2)=I(1:2)+I(2:1)\).

(4) Apply (E2) to \(u=f_1\log(f_1/f_2)\), whose negative part is
integrable by (1). If \(\mu_1(E_j)>0\), then
\(\int_{E_j}u\,d\lambda=\mu_1(E_j)I_{1:2}(E_j)\). If \(\mu_1(E_j)=0\),
then \(\lambda(E_j)=0\), so \(\int_{E_j}u\,d\lambda=0\); this agrees with
the convention \(I_{1:2}(E_j)=0\).
\end{proof}
\end{annexitem}

\begin{annexitem}{annex:p15-w02b-p15-045-extension}{thm:p15-w02b-p15-045-extension}{The extended values are well defined and agree with the source's}
\begin{proof}
Put \(\psi(t)=\varphi(t)-t+1\) for \(t\geq0\), so \(\psi(0)=1\).
\(\varphi\) is continuous on \([0,\infty)\), since \(t\log t\to0\) as
\(t\downarrow0\). On \((0,\infty)\), \(\psi'(t)=\log t\) is negative on
\((0,1)\) and positive on \((1,\infty)\), and \(\psi(1)=0\). With
continuity at \(0\), this gives \(\psi(t)\geq0\) for all \(t\geq0\),
with equality only at \(t=1\). Also \(\varphi\geq-1/e\) on
\([0,\infty)\), by part~(1) of the proof of
\cref{thm:p15-w02b-p15-045-well-defined} and \(\varphi(0)=0\).

(1) Let \(Z=\{g=0\}\). Then \(\mu_1(Z)=\int_Zg\,d\mu_2=0\), so
\(\log g\) is real \(\mu_1\)-almost everywhere. The change-of-measure
rule (\cref{thm:p15-density-calculus}, part~1, with \(\nu=\mu_1\),
\(\mu=\mu_2\), density \(g\)) applied to the nonnegative functions
\(1_{X\setminus Z}(\log g)^\pm\) gives
\[
  \int(\log g)^\pm\,d\mu_1=\int_{X\setminus Z}g(\log g)^\pm\,d\mu_2
  =\int\varphi(g)^\pm\,d\mu_2,
\]
because \(g(\log g)^\pm=(g\log g)^\pm\) where \(g>0\), and
\(\varphi(g)=0\) on \(Z\). The negative part is at most
\(\mu_2(X)/e=1/e\). So the extended integral exists and
\(I(1:2)=\int\varphi(g)\,d\mu_2\). Since \(\mu_1\ll\mu_2\),
\(\int g\,d\mu_2=\mu_1(X)=1\), and \(g-1\) is \(\mu_2\)-integrable with
integral \(0\). Pointwise \(\varphi(g)=(g-1)+\psi(g)\), so (E1) gives
\(I(1:2)=\int\psi(g)\,d\mu_2\in[0,\infty]\). It is \(0\) if and only if
\(\psi(g)=0\) \([\mu_2]\), that is, \(g=1\) \([\mu_2]\). If \(g=1\)
\([\mu_2]\), then \(\mu_1(A)=\int_Ag\,d\mu_2=\mu_2(A)\) for every
\(A\). Conversely, if \(\mu_1=\mu_2\), the constant \(1\) is a version of
\(d\mu_1/d\mu_2\), and two versions agree \([\mu_2]\)
(\cref{thm:p15-radon-nikodym}). Two versions of \(g\) agree
\([\mu_2]\), hence \([\mu_1]\), so \(\int\log g\,d\mu_1\) does not
depend on the version.

(2) Under the standing hypotheses, \(f_1/f_2\) is a version of
\(d\mu_1/d\mu_2\) with values in \((0,\infty)\) \([\lambda]\)
(\cref{thm:p15-w02b-p15-045-well-defined}(2)), and part~(2) of that
proof gives
\[
  \int f_1\log\frac{f_1}{f_2}\,d\lambda=\int\log\frac{f_1}{f_2}\,d\mu_1 .
\]
The left side is \(I(1:2)\) of \cref{def:p15-w02b-p15-045-information},
and the right side is \(I(1:2)\) of
\cref{def:p15-w02b-p15-045-extension}(1) computed with the version
\(f_1/f_2\); by (1) the version does not matter. Exchanging the indices
gives \(I(2:1)\). The two definitions of \(J(1,2)\) then agree by
\cref{thm:p15-w02b-p15-045-well-defined}(3).

(3) If \(\mu_1\ll\mu_2\) and \(\mu_2\) is not absolutely continuous with
respect to \(\mu_1\), then \(\mu_1\neq\mu_2\), so \(I(1:2)>0\) by (1),
and \(I(2:1)=+\infty\) by clause (2) of the definition with the indices
exchanged. Then \(J(1,2)=I(1:2)+\infty=+\infty\). If neither measure is
absolutely continuous with respect to the other, clause (2) applies in
both directions.
\end{proof}
\end{annexitem}

\begin{annexitem}{annex:p15-w02b-p15-045-nonnegative}{thm:p15-w02b-p15-045-nonnegative}{Lemma 3.1: almost positive definiteness}
\begin{proof}
Let \(g=f_1/f_2\), finite and positive \([\lambda]\). By
\cref{thm:p15-w02b-p15-045-well-defined}(2), \(g=d\mu_1/d\mu_2\) and
\(I(1:2)=\int\log g\,d\mu_1\). The change-of-measure rule applied to
\((\log g)^\pm\), with \(\nu=\mu_1\), \(\mu=\mu_2\) and density \(g\),
gives \(\int(\log g)^\pm\,d\mu_1=\int g(\log g)^\pm\,d\mu_2\). Since
\(g>0\), \(g(\log g)^\pm=(g\log g)^\pm\). Hence
\[
  I(1:2)=\int g\log g\,d\mu_2,\qquad \int g\,d\mu_2=\mu_1(X)=1,
\]
which are the source's (3.1) and (3.2).

For \(t>0\) put \(\psi(t)=t\log t-t+1\). Then \(\psi(1)=0\) and
\(\psi'(t)=\log t\), which is negative on \((0,1)\) and positive on
\((1,\infty)\). Hence \(\psi(t)\geq0\), with equality only at \(t=1\).
This is the source's second-order expansion (3.3): with
\(\varphi(t)=t\log t\), Taylor's theorem gives
\(\varphi(t)=\varphi(1)+(t-1)\varphi'(1)+\tfrac12(t-1)^2\varphi''(h)\)
for some \(h\) between \(t\) and \(1\). Since \(\varphi(1)=0\),
\(\varphi'(1)=1\) and \(\varphi''(h)=1/h>0\), this says
\(\psi(t)=\tfrac12(t-1)^2/h\geq0\).

Now \(g\log g=(g-1)+\psi(g)\) \([\lambda]\). The function \(g-1\) is
\(\mu_2\)-integrable with integral \(0\), and \(\psi(g)\geq0\). By (E1),
\[
  I(1:2)=\int(g-1)\,d\mu_2+\int\psi(g)\,d\mu_2=\int\psi(g)\,d\mu_2\in[0,\infty].
\]
This is (3.4)--(3.5), and it holds whether or not \(I(1:2)\) is finite.
The value is \(0\) if and only if \(\psi(g)=0\) \(\mu_2\)-almost
everywhere, that is, \(g=1\) \([\mu_2]\), equivalently \(f_1=f_2\)
\([\lambda]\) because \(\mu_2\equiv\lambda\). If \(f_1=f_2\) \([\lambda]\),
the two measures have the same integrals over every set, so
\(\mu_1=\mu_2\). Conversely, \(\mu_1=\mu_2\) makes \(f_1,f_2\) two
versions of one density, which agree \([\lambda]\)
(\cref{thm:p15-radon-nikodym}). The hypotheses are symmetric in the
indices, so the same argument applies to \(I(2:1)\). Finally,
\(J(1,2)=I(1:2)+I(2:1)\) by
\cref{thm:p15-w02b-p15-045-well-defined}(3). It is a sum of two elements
of \([0,\infty]\), and it vanishes exactly when both terms do, that is,
when \(\mu_1=\mu_2\).
\end{proof}
\end{annexitem}

\begin{annexitem}{annex:p15-w02b-p15-045-event-bound}{thm:p15-w02b-p15-045-event-bound}{Lemma 3.2: information from an event}
\begin{proof}
Since \(\lambda(E)>0\) and \(\lambda\equiv\{\mu_1,\mu_2\}\), both
\(\mu_1(E)\) and \(\mu_2(E)\) are positive, and
\(c=\log(\mu_1(E)/\mu_2(E))\) is a real number. If
\(I_{1:2}(E)=\infty\), then \(I_{1:2}(E)>c\). Otherwise, and in fact
in general, consider the measurable space \((E,\mathbf S\cap E)\) with
\[
  \mu_i^E(A)=\frac{\mu_i(A)}{\mu_i(E)},\qquad
  \lambda^E(A)=\frac{\lambda(A)}{\lambda(E)},\qquad
  f_i^E=\frac{\lambda(E)}{\mu_i(E)}\,f_i\Big|_E .
\]
These are probability measures. \(\int_Af_i^E\,d\lambda^E=\mu_i(A)/\mu_i(E)=\mu_i^E(A)\),
\(0<f_i^E<\infty\) \([\lambda^E]\), and the three measures have the same
null sets as \(\lambda\) on \(E\). So the standing hypotheses hold on
\(E\). The information of this pair is
\[
  I^E(1:2)=\int_Ef_1^E\log\frac{f_1^E}{f_2^E}\,d\lambda^E
          =\frac1{\mu_1(E)}\int_Ef_1\Bigl(\log\frac{f_1}{f_2}-c\Bigr)d\lambda
          =I_{1:2}(E)-c,
\]
where the last step is (E1) applied to \(u=1_Ef_1\log(f_1/f_2)/\mu_1(E)\)
(integrable negative part, \cref{thm:p15-w02b-p15-045-well-defined}(1))
and the integrable function \(v=-c\,1_Ef_1/\mu_1(E)\), whose integral is
\(-c\). This is the source's displayed identity in the proof of
Lemma~3.2. By \cref{thm:p15-w02b-p15-045-nonnegative},
\(I^E(1:2)\geq0\), that is, \(I_{1:2}(E)\geq c\). When
\(I_{1:2}(E)<\infty\), equality holds if and only if \(I^E(1:2)=0\). By
the same lemma this happens if and only if \(f_1^E=f_2^E\)
\([\lambda^E]\), which says \(f_1/f_2=\mu_1(E)/\mu_2(E)\) for
\(\lambda\)-almost every \(x\in E\).
\end{proof}
\end{annexitem}

\begin{annexitem}{annex:p15-w02b-p15-045-additivity}{thm:p15-w02b-p15-045-additivity}{Theorem 3.1: additivity for independent components}
\begin{proof}
\emph{Densities.} \(f_i(x,y)=f_i^{(1)}(x)f_i^{(2)}(y)\) is
\(\mathbf S\)-measurable, being the product of the measurable functions
\(f_i^{(1)}\) of the first coordinate and \(f_i^{(2)}\) of the second. For a rectangle
\(A\times B\), Tonelli's theorem (\cref{thm:p15-tonelli}) gives
\[
  \int_{A\times B}f_i\,d\lambda
   =\int_Af_i^{(1)}\,d\lambda_1\int_Bf_i^{(2)}\,d\lambda_2
   =\mu_i^{(1)}(A)\,\mu_i^{(2)}(B)=\mu_i(A\times B).
\]
The measure \(E\mapsto\int_Ef_i\,d\lambda\) is finite and agrees with
\(\mu_i\) on rectangles. By the uniqueness part of
\cref{thm:p15-product-existence} it equals \(\mu_i\). Let \(N_k\) be the
\(\lambda_k\)-null set on which \(f_1^{(k)}\) or \(f_2^{(k)}\) lies
outside \((0,\infty)\). Then \(f_i\in(0,\infty)\) off
\((N_1\times X_2)\cup(X_1\times N_2)\), which is \(\lambda\)-null.
\(\lambda\) is a probability measure. If \(\lambda(E)=0\), then
\(\mu_i(E)=\int_Ef_i\,d\lambda=0\). If \(\mu_i(E)=0\), then
\(\int_Ef_i\,d\lambda=0\) with \(f_i>0\) \([\lambda]\) forces
\(\lambda(E)=0\). So \(\lambda\equiv\{\mu_1,\mu_2\}\), \(\mu_1\equiv\mu_2\),
and the standing hypotheses hold.

\emph{Additivity.} Off a \(\lambda\)-null set,
\(\log(f_1/f_2)(x,y)=a(x)+b(y)\) with
\(a=\log(f_1^{(1)}/f_2^{(1)})\) and \(b=\log(f_1^{(2)}/f_2^{(2)})\).
Hence \(f_1\log(f_1/f_2)=u+v\), where
\(u(x,y)=\bigl(f_1^{(1)}a\bigr)(x)\,f_1^{(2)}(y)\) and
\(v(x,y)=f_1^{(1)}(x)\,\bigl(f_1^{(2)}b\bigr)(y)\). Since
\(f_1^{(2)}>0\), \(u^\pm(x,y)=(f_1^{(1)}a)^\pm(x)f_1^{(2)}(y)\), and
Tonelli gives
\[
  \int u^\pm\,d\lambda=\int(f_1^{(1)}a)^\pm\,d\lambda_1\cdot\int f_1^{(2)}\,d\lambda_2
                     =\int(f_1^{(1)}a)^\pm\,d\lambda_1 .
\]
By \cref{thm:p15-w02b-p15-045-well-defined}(1) for the first factor,
\(\int u^-\,d\lambda<\infty\) and \(\int u\,d\lambda=I_x(1:2)\). In the
same way \(\int v^-\,d\lambda<\infty\) and \(\int v\,d\lambda=I_y(1:2)\).
By (E1), \(I_{xy}(1:2)=\int(u+v)\,d\lambda=I_x(1:2)+I_y(1:2)\) in
\([0,\infty]\). This is the source's computation (3.6), with the
splitting justified.

Exchanging the indices gives the statement for \(I(2:1)\). Adding the
two statements and using \(J=I(1:2)+I(2:1)\)
(\cref{thm:p15-w02b-p15-045-well-defined}(3)) gives it for \(J\). For
\(n\) observations, the \(n\)-fold product of \(\mu_i\) is the product of
the \((n-1)\)-fold product with one more factor
(\cref{thm:p15-finite-product-measures}). Induction on \(n\) with the
two-factor case gives \(n\,I(1:2)\).
\end{proof}
\end{annexitem}

\begin{annexitem}{annex:p15-w02b-p15-045-induced}{thm:p15-w02b-p15-045-induced}{The induced pair and the conditional-expectation form of its densities}
\begin{proof}
(1) \(\gamma,\nu_1,\nu_2\) are probability measures on \(\mathbf T\)
(\cref{thm:p15-pushforward-laws}). For \(G\in\mathbf T\), the following
are equivalent: \(\gamma(G)=0\); \(\lambda(T^{-1}G)=0\);
\(\mu_1(T^{-1}G)=0\); \(\mu_2(T^{-1}G)=0\); \(\nu_1(G)=0\); \(\nu_2(G)=0\).
So \(\nu_1\equiv\nu_2\) and \(\gamma\equiv\{\nu_1,\nu_2\}\). By
\cref{thm:p15-radon-nikodym}, \(g_i=d\nu_i/d\gamma\) exists and is finite
\([\gamma]\). Also
\(\nu_i(\{g_i=0\})=\int_{\{g_i=0\}}g_i\,d\gamma=0\), so
\(\gamma(\{g_i=0\})=0\). Redefining \(g_i\) to be \(1\) on the
measurable \(\gamma\)-null set where it lies outside \((0,\infty)\)
produces a version with \(0<g_i<\infty\) everywhere.

(2) \(g_i\circ T\) is \(T^{-1}(\mathbf T)\)-measurable, since
\(\{g_i\circ T\in B\}=T^{-1}(g_i^{-1}B)\). \emph{Transfer:} for every
\(\mathbf T\)-measurable \(u\geq0\),
\(\int_Xu\circ T\,d\lambda=\int_Yu\,d\gamma\) and
\(\int_Xu\circ T\,d\mu_i=\int_Yu\,d\nu_i\). For \(u=1_G\) this is the
definition of \(\gamma\) and \(\nu_i\). It extends to simple \(u\) by
linearity, and to general \(u\) by monotone convergence along simple
\(u_n\uparrow u\), since then \(u_n\circ T\uparrow u\circ T\). Hence
\[
  \int_{T^{-1}G}g_i\circ T\,d\lambda=\int_X(1_Gg_i)\circ T\,d\lambda
   =\int_Gg_i\,d\gamma=\nu_i(G),
\]
and \(\nu_i(G)=\mu_i(T^{-1}G)=\int_{T^{-1}G}f_i\,d\lambda\).
With \(G=Y\) this shows that \(g_i\circ T\) is integrable. Every
member of \(T^{-1}(\mathbf T)\) is of the form \(T^{-1}G\). So
\(g_i\circ T\) satisfies the two defining requirements of
\cref{def:p15-general-conditional-expectation} for the conditional
expectation of \(f_i\) given \(T^{-1}(\mathbf T)\) under \(\lambda\).

(3) Let \(\nu_1(G)>0\). By (1),
\cref{thm:p15-w02b-p15-045-well-defined}(1) applies on
\((Y,\mathbf T)\). So the integral of \(g_1\log(g_1/g_2)\) over \(G\)
with respect to \(\gamma\) exists in \((-\infty,\infty]\). The change-of-measure rule applied to the positive
and negative parts, with density \(g_1=d\nu_1/d\gamma\), gives
\(\int_Gg_1\log(g_1/g_2)\,d\gamma=\int_G\log(g_1/g_2)\,d\nu_1\). Transfer
with \((\mu_1,\nu_1)\), applied to \((1_G\log(g_1/g_2))^\pm\), turns this
into \(\int_{T^{-1}G}\log(g_1T/g_2T)\,d\mu_1\). Dividing by
\(\nu_1(G)=\mu_1(T^{-1}G)\) gives the first formula, and \(G=Y\) gives
the second.
\end{proof}
\end{annexitem}

\begin{annexitem}{annex:p15-w02b-p15-045-density-criterion}{thm:p15-w02b-p15-045-density-criterion}{The density step: sufficiency as measurability of the ratio}
\begin{proof}
Write \(\mathbf S_0=T^{-1}(\mathbf T)\). Let \(h\) be a version of
\(f_1/f_2\) with values in \((0,\infty)\) everywhere; this changes
\(f_1/f_2\) only on a \(\lambda\)-null set. Put \(h'=g_1T/g_2T\), which is
\(\mathbf S_0\)-measurable with values in \((0,\infty)\). By
\cref{thm:p15-w02b-p15-045-well-defined}(2), \(h=d\mu_1/d\mu_2\). An
\(\mathbf S_0\)-measurable function has the same integral with respect to
a measure on \(\mathbf S\) as with respect to its restriction to
\(\mathbf S_0\), since its simple approximations are
\(\mathbf S_0\)-measurable.

(3)\(\Rightarrow\)(2): take \(W=h'\).

(2)\(\Rightarrow\)(3): let \(h=W\) \([\lambda]\) with \(W\)
\(\mathbf S_0\)-measurable. We may take \(W\geq0\), modifying \(W\) on the
\(\mathbf S_0\)-measurable \(\lambda\)-null set \(\{W<0\}\). Then
\(f_1=Wf_2\) \([\lambda]\). For \(A\in\mathbf S_0\),
\cref{thm:p15-w02b-p15-045-induced}(2) and (P), with \(\rho=\lambda\),
\(V=f_2\), \(V_0=g_2T\) and multiplier \(1_AW\), give
\[
  \int_Ag_1T\,d\lambda=\int_Af_1\,d\lambda=\int1_AW\,f_2\,d\lambda
  =\int1_AW\,g_2T\,d\lambda .
\]
In particular \(Wg_2T\) is integrable. So \(g_1T\) and \(Wg_2T\) are
integrable \(\mathbf S_0\)-measurable functions with the same integral
over every \(A\in\mathbf S_0\). On \(A=\{g_1T>Wg_2T\}\in\mathbf S_0\) the
integral of the positive difference is \(0\), so \(\lambda(A)=0\), and
likewise for the reverse inequality. Hence \(g_1T=Wg_2T\) \([\lambda]\),
that is, \(h'=W=h\) \([\lambda]\).

(2)\(\Rightarrow\)(1): with \(W\geq0\) as above, \(W=h=d\mu_1/d\mu_2\)
\([\mu_2]\). For \(E\in\mathbf S\), let \(\phi_E\) be a version of the
conditional expectation of \(1_E\) given \(\mathbf S_0\) under \(\mu_2\)
(\cref{thm:p15-conditional-existence}). Taking
\(A=\{\phi_E<0\}\in\mathbf S_0\) in
\(\int_A\phi_E\,d\mu_2=\mu_2(E\cap A)\geq0\) shows \(\mu_2(A)=0\). So
\(\phi_E\) may be taken nonnegative. Then
\(\mu_2(E\cap A)=\int_A\phi_E\,d\mu_2\) for \(A\in\mathbf S_0\) by
definition. For \(\mu_1\), the change-of-measure rule, then (P) with
\(\rho=\mu_2\), \(V=1_E\), \(V_0=\phi_E\) and multiplier \(1_AW\), and
then the change-of-measure rule again give
\[
  \mu_1(E\cap A)=\int1_E1_AW\,d\mu_2=\int\phi_E1_AW\,d\mu_2=\int_A\phi_E\,d\mu_1 .
\]
So one \(\mathbf S_0\)-measurable \(\phi_E\) serves both measures, and
\(T\) is sufficient.

(1)\(\Rightarrow\)(2): let \(T\) be sufficient with functions
\(\phi_E\). As in the previous step, \(\phi_E\geq0\)
\([\mu_2]\); modify it on an \(\mathbf S_0\)-measurable null set to be
nonnegative everywhere. The restriction of \(\mu_1\) to \(\mathbf S_0\)
is absolutely continuous with respect to that of \(\mu_2\). Let \(W\) be
the \(\mathbf S_0\)-measurable density
(\cref{thm:p15-radon-nikodym} on \((X,\mathbf S_0)\)), taken
nonnegative and finite. For \(E\in\mathbf S\), sufficiency with
\(A=X\), the change-of-measure rule on \(\mathbf S_0\) for the
\(\mathbf S_0\)-measurable \(\phi_E\geq0\), and (P) with \(\rho=\mu_2\),
\(V=1_E\), \(V_0=\phi_E\) and multiplier \(W\) give
\[
  \mu_1(E)=\int\phi_E\,d\mu_1=\int\phi_EW\,d\mu_2=\int_EW\,d\mu_2 .
\]
So \(W\) is a version of \(d\mu_1/d\mu_2\) on \(\mathbf S\). Hence
\(W=h\) \([\mu_2]\), that is, \([\lambda]\), and (2) holds.
\end{proof}
\end{annexitem}

\begin{annexitem}{annex:p15-w02b-p15-045-reduction}{thm:p15-w02b-p15-045-reduction}{Theorem 4.1: data reduction, with the exact finiteness condition}
\begin{proof}
Use \(\mathbf S_0\), \(h\) and \(h'\) as in the proof of
\cref{thm:p15-w02b-p15-045-density-criterion}, with \(0<g_i<\infty\)
everywhere (\cref{thm:p15-w02b-p15-045-induced}(1)). By
\cref{thm:p15-w02b-p15-045-well-defined}(2),
\(I(1:2)=\int\log h\,d\mu_1\). By
\cref{thm:p15-w02b-p15-045-induced}(3),
\(I'(1:2)=\int\log h'\,d\mu_1\). Both lie in \([0,\infty]\) by
\cref{thm:p15-w02b-p15-045-nonnegative}, applied on \(X\) and on \(Y\).

\emph{The comparison measure.} Define
\(\mu_{12}(E)=\int_Ef_2h'\,d\lambda\) for \(E\in\mathbf S\), as on p.~84
of the source. By (P) with \(\rho=\lambda\), \(V=f_2\), \(V_0=g_2T\)
(\cref{thm:p15-w02b-p15-045-induced}(2)) and multiplier \(h'\), and then
by transfer,
\[
  \mu_{12}(X)=\int h'g_2T\,d\lambda=\int g_1T\,d\lambda=\int g_1\,d\gamma=1 .
\]
So \(\mu_{12}\) is a probability measure. Its density \(f_2h'\) lies in
\((0,\infty)\) \([\lambda]\), so \(\mu_{12}\equiv\lambda\). Put
\(g=h/h'=f_1g_2T/(f_2g_1T)\in(0,\infty)\) \([\lambda]\). By the
change-of-measure rule,
\(\int g\,d\mu_{12}=\int(h/h')f_2h'\,d\lambda=\int f_1\,d\lambda=1\).

\emph{The gap.} The proof of \cref{thm:p15-w02b-p15-045-nonnegative}
used only that \(g\) is finite and positive almost everywhere and
\(\int g\,d\mu_2=1\). Repeated with \(\mu_{12}\) in place of \(\mu_2\),
it shows that
\(D=\int g\log g\,d\mu_{12}=\int\psi(g)\,d\mu_{12}\) exists in
\([0,\infty]\), with \(D=0\) if and only if \(g=1\) \([\mu_{12}]\), that
is, \(h=h'\) \([\lambda]\). The change-of-measure rule on positive and
negative parts gives \(D=\int f_1\log g\,d\lambda\), because
\(g\cdot f_2h'=f_1\). This is (4.6).

\emph{The identity \(I(1:2)=I'(1:2)+D\).} Off a \(\lambda\)-null set,
\(f_1\log h=f_1\log h'+f_1\log g\). The negative part of \(f_1\log h'\)
is \(\lambda\)-integrable, because \(\int(\log h')^-\,d\mu_1<\infty\) by
\cref{thm:p15-w02b-p15-045-induced}(3) and
\cref{thm:p15-w02b-p15-045-well-defined}(1) on \(Y\). The negative part
of \(f_1\log g\) is integrable, because
\(\int(f_1\log g)^-\,d\lambda=\int(g\log g)^-\,d\mu_{12}\leq1/e\). By
(E1), in \([0,\infty]\),
\[
  I(1:2)=I'(1:2)+D .
\]
This replaces the source's subtraction (4.5), which is undefined when
both terms are infinite, by an addition that is always defined.

(1) Since \(D\geq0\), \(I(1:2)\geq I'(1:2)\).

(2) If \(T\) is sufficient, \cref{thm:p15-w02b-p15-045-density-criterion}
gives \(h=h'\) \([\lambda]\). Hence \(g=1\), \(D=0\) and
\(I(1:2)=I'(1:2)\), finite or not.

(3) If \(I(1:2)<\infty\), then \(I'(1:2)\leq I(1:2)<\infty\) and
\(D=I(1:2)-I'(1:2)\) is finite. Equality \(I(1:2)=I'(1:2)\) holds if and
only if \(D=0\), if and only if \(h=h'\) \([\lambda]\) (the source's
(4.7)), and by \cref{thm:p15-w02b-p15-045-density-criterion} if and only
if \(T\) is sufficient.

(4) This is \cref{thm:p15-w02b-p15-045-infinite-equality}(1), proved in
\cref{annex:p15-w02b-p15-045-infinite-equality}. There
\(I(1:2)=I'(1:2)=\infty\), and \(T\) is not sufficient. The identity
\(I=I'+D\) then reads \(\infty=\infty+D\) with \(D>0\), so equality of
the two informations carries no information about \(D\).
\end{proof}
\end{annexitem}

\begin{annexitem}{annex:p15-w02b-p15-045-nonsingular}{thm:p15-w02b-p15-045-nonsingular}{Corollary 4.1 and the divergence form}
\begin{proof}
(1) If \(T^{-1}(\mathbf T)=\mathbf S\), then \(f_1/f_2\) is
\(T^{-1}(\mathbf T)\)-measurable. Condition (2) of
\cref{thm:p15-w02b-p15-045-density-criterion} holds, so \(T\) is
sufficient, and \cref{thm:p15-w02b-p15-045-reduction}(2) gives
\(I(1:2)=I'(1:2)\).

(2) By \cref{thm:p15-w02b-p15-045-well-defined}(3) on \(X\), and on
\(Y\) through \cref{thm:p15-w02b-p15-045-induced}(1),
\(J=I(1:2)+I(2:1)\) and \(J'=I'(1:2)+I'(2:1)\). The standing hypotheses,
the construction of \(\nu_i,\gamma,g_i\), and the definition of
sufficiency are symmetric in the indices.
\cref{thm:p15-w02b-p15-045-reduction} therefore also gives
\(I(2:1)\geq I'(2:1)\), with equality when \(T\) is sufficient. Adding
the two inequalities gives \(J\geq J'\), with equality when \(T\) is
sufficient. If \(J<\infty\), then \(I(1:2)\) and \(I(2:1)\) are finite,
and \(J=J'\) says that the two nonnegative finite gaps
\(I(1:2)-I'(1:2)\) and \(I(2:1)-I'(2:1)\) sum to \(0\). So both vanish,
and \cref{thm:p15-w02b-p15-045-reduction}(3) gives sufficiency.
\end{proof}
\end{annexitem}

\begin{annexitem}{annex:p15-w02b-p15-045-events}{thm:p15-w02b-p15-045-events}{Theorem 4.2}
\begin{proof}
Let \(G\in\mathbf T\). If \(\nu_1(G)=0\), then \(\mu_1(T^{-1}G)=0\) and
both quantities are \(0\) by convention. If \(\nu_1(G)>0\), then
\(\mu_1(T^{-1}G)=\nu_1(G)>0\). By
\cref{thm:p15-w02b-p15-045-well-defined}(2),
\(I_{1:2}(T^{-1}G)=\mu_1(T^{-1}G)^{-1}\int_{T^{-1}G}\log h\,d\mu_1\). By
\cref{thm:p15-w02b-p15-045-induced}(3),
\(I'_{1:2}(G)=\mu_1(T^{-1}G)^{-1}\int_{T^{-1}G}\log h'\,d\mu_1\). This is
(4.9).

(1) Take \(G=Y\); then \(T^{-1}G=X\).

(2) If \(T\) is sufficient, \(h=h'\) \([\lambda]\), hence \([\mu_1]\)
(\cref{thm:p15-w02b-p15-045-density-criterion}). The two integrals then
agree for every \(G\).

(3) Suppose \(I(1:2)<\infty\). If \(I(1:2)=I'(1:2)\),
\cref{thm:p15-w02b-p15-045-reduction}(3) makes \(T\) sufficient, and
(2) gives equality for every \(G\). The converse is (1).

The final assertion is \cref{thm:p15-w02b-p15-045-infinite-equality}(2).
\end{proof}
\end{annexitem}

\begin{annexitem}{annex:p15-w02b-p15-045-j-counterparts}{thm:p15-w02b-p15-045-j-counterparts}{Section 5: the divergence counterparts of Lemma 3.2, Corollary 4.1 and Theorem 4.2}
\begin{proof}
The standing hypotheses, the construction of \(\nu_i,\gamma,g_i\) in
\cref{def:p15-w02b-p15-045-statistic} and the definition of sufficiency
are symmetric in the indices \(1\) and \(2\). Each result about
\(I_{1:2}\), \(I(1:2)\) and \(I'(1:2)\) used below therefore also holds
for \(I_{2:1}\), \(I(2:1)\) and \(I'(2:1)\). By definition,
\(J_{12}(E)=I_{1:2}(E)+I_{2:1}(E)\) and
\(J'_{12}(G)=I'_{1:2}(G)+I'_{2:1}(G)\), sums of elements of
\((-\infty,\infty]\) (\cref{thm:p15-w02b-p15-045-well-defined}(1), on
\(X\) and, through \cref{thm:p15-w02b-p15-045-induced}(1), on \(Y\)).

(1) Since \(\lambda(E)>0\) and \(\lambda\equiv\{\mu_1,\mu_2\}\), both
\(\mu_1(E)\) and \(\mu_2(E)\) are positive; put
\(c=\log(\mu_1(E)/\mu_2(E))\in\mathbb R\).
\Cref{thm:p15-w02b-p15-045-event-bound} gives \(I_{1:2}(E)\geq c\), and
the same lemma with the indices exchanged gives \(I_{2:1}(E)\geq-c\).
Addition on \((-\infty,\infty]\) is monotone in each argument, so
\(J_{12}(E)\geq c+(-c)=0\). If \(J_{12}(E)=\infty\), then
\(J_{12}(E)>0\). Suppose \(J_{12}(E)<\infty\). Neither term is
\(-\infty\), so both are finite, and
\[
  J_{12}(E)=\bigl(I_{1:2}(E)-c\bigr)+\bigl(I_{2:1}(E)+c\bigr),
\]
a sum of two nonnegative real numbers. It is \(0\) if and only if both
are \(0\). By the equality clause of
\cref{thm:p15-w02b-p15-045-event-bound}, the first vanishes if and only
if \(f_1/f_2=\mu_1(E)/\mu_2(E)\) for \(\lambda\)-almost every \(x\in E\),
and the second, by the same clause with the indices exchanged, if and
only if \(f_2/f_1=\mu_2(E)/\mu_1(E)\) for \(\lambda\)-almost every
\(x\in E\). Because \(0<f_i<\infty\) \([\lambda]\), the two conditions
are the same. So \(J_{12}(E)=0\) if and only if that condition holds.

(2) If \(T^{-1}(\mathbf T)=\mathbf S\), then \(f_1/f_2\) is
\(T^{-1}(\mathbf T)\)-measurable, so \(T\) is sufficient by
\cref{thm:p15-w02b-p15-045-density-criterion}, as in the proof of
\cref{thm:p15-w02b-p15-045-nonsingular}(1). By
\cref{thm:p15-w02b-p15-045-nonsingular}(2), sufficiency gives
\(J(1,2)=J'(1,2)\), whether or not the value is finite.

(3)(a) Take \(G=Y\). Then \(T^{-1}G=X\), \(J_{12}(X)=J(1,2)\) and
\(J'_{12}(Y)=J'(1,2)\).

(b) If \(T\) is sufficient, \cref{thm:p15-w02b-p15-045-events}(2) gives
\(I_{1:2}(T^{-1}G)=I'_{1:2}(G)\) for every \(G\in\mathbf T\), and the
same result with the indices exchanged gives
\(I_{2:1}(T^{-1}G)=I'_{2:1}(G)\). Adding the two equalities gives
\(J_{12}(T^{-1}G)=J'_{12}(G)\).

(c) Suppose \(J(1,2)<\infty\). By
\cref{thm:p15-w02b-p15-045-nonsingular}(2), \(J(1,2)=J'(1,2)\) if and only
if \(T\) is sufficient. If \(T\) is sufficient, (b) gives the equalities
for every \(G\). If the equalities hold for every \(G\), (a) gives
\(J(1,2)=J'(1,2)\). So the three conditions are equivalent.

The final assertion is \cref{thm:p15-w02b-p15-045-infinite-equality}(1)
and~(3): there \(J(1,2)=J'(1,2)=\infty\), \(T\) is not sufficient, and
for \(G=(\tfrac12,1)\) the event equality fails.
\end{proof}
\end{annexitem}

\begin{annexitem}{annex:p15-w02b-p15-045-fisher}{thm:p15-w02b-p15-045-fisher}{Equations (2.7)--(2.9): the second-order relation to Fisher's information}
\begin{proof}
Let \(B^\circ=\{\theta':|\theta'-\theta|<r\}\). All partial derivatives
below are taken in the parameter, for fixed \(x\).

\emph{Measurability.} For fixed \(\theta'\) in the open set of (R1),
\(\partial_\alpha f(\cdot,\theta')\) is the pointwise limit of the
\(\mathbf S\)-measurable difference quotients
\(n\bigl(f(\cdot,\theta'+e_\alpha/n)-f(\cdot,\theta')\bigr)\), with
\(e_\alpha\) the \(\alpha\)-th unit vector, so it is
\(\mathbf S\)-measurable. The same argument applies to second
derivatives and to the derivatives of \(\ell=\log f\), which are finite
because \(f>0\).

\emph{Step 1: differentiation under the integral.} Let \(\phi(x,\theta')\)
be continuously differentiable in \(\theta'\) on an open set containing
\(B\), with \(\int|\phi(\cdot,\theta')|\,d\lambda<\infty\) for
\(\theta'\in B^\circ\) and \(|\partial_\beta\phi(x,\theta')|\leq N(x)\)
for \(\theta'\in B\). Fix \(\theta'\in B^\circ\) and a sequence
\(t_n\to0\), \(t_n\neq0\), with \(\theta'+t_ne_\beta\in B^\circ\). For each
\(x\) the mean value theorem in the single variable \(t\) gives
\(\bigl(\phi(x,\theta'+t_ne_\beta)-\phi(x,\theta')\bigr)/t_n=\partial_\beta\phi(x,\theta'+\tau e_\beta)\)
for some \(\tau\) between \(0\) and \(t_n\). So the quotient is bounded
by \(N(x)\), and it tends to \(\partial_\beta\phi(x,\theta')\) by
continuity. Dominated convergence
(\cref{thm:p15-dominated-bounded-convergence}) gives
\[
  \frac{\int\phi(\cdot,\theta'+t_ne_\beta)\,d\lambda-\int\phi(\cdot,\theta')\,d\lambda}{t_n}
  \to\int\partial_\beta\phi(\cdot,\theta')\,d\lambda .
\]
With \(\phi=f\), the left side is \(0\) because
\(\int f(\cdot,\theta'')\,d\lambda=1\) for every \(\theta''\), and (R2)
supplies the bound. Hence
\(\int\partial_\beta f(\cdot,\theta')\,d\lambda=0\) for
\(\theta'\in B^\circ\). With \(\phi=\partial_\alpha f\), integrable on
\(B^\circ\) by (R2), the left side is again \(0\) by what was just
shown, and (R2) bounds \(\partial_\beta\partial_\alpha f\). Hence
\(\int\partial_\beta\partial_\alpha f(\cdot,\theta')\,d\lambda=0\) for
\(\theta'\in B^\circ\).

\emph{Step 2: part (1).} At \((x,\theta)\),
\(\partial_\alpha\ell=\partial_\alpha f/f\) and
\(\partial_\alpha\partial_\beta\ell=\partial_\alpha\partial_\beta f/f-(\partial_\alpha f)(\partial_\beta f)/f^2\).
Hence
\[
  f\Bigl(\frac1f\partial_\alpha f\Bigr)\Bigl(\frac1f\partial_\beta f\Bigr)
  =\partial_\alpha\partial_\beta f-f\,\partial_\alpha\partial_\beta\ell .
\]
The first term on the right is bounded by \(N\) and the second by \(K\)
((R2), (R3) with \(\theta'=\theta''=\theta\)), so the left side is
\(\lambda\)-integrable, \(g_{\alpha\beta}\) is finite, and by Step~1
\(g_{\alpha\beta}=0-\int f\,\partial_\alpha\partial_\beta\ell\,d\lambda\).
Symmetry is clear from the defining integral. For \(v\in\mathbb R^k\),
linearity gives
\(\sum g_{\alpha\beta}v_\alpha v_\beta=\int f\bigl(\sum_\alpha v_\alpha\partial_\alpha\ell\bigr)^2d\lambda\geq0\).

\emph{Step 3: Taylor's formula.} Fix \(\Delta\theta\) with
\(0<|\Delta\theta|<r\), write \(\Delta=\Delta\theta\), and put
\(\|\Delta\|_1=\sum_\alpha|\Delta_\alpha|\), so
\(\|\Delta\|_1^2\leq k|\Delta|^2\) by the Cauchy--Schwarz inequality.
The segment from \(\theta\) to \(\theta+\Delta\) lies in \(B^\circ\).
For fixed \(x\) and a base point \(\eta\in\{\theta,\theta+\Delta\}\)
with step \(D=\pm\Delta\) such that \(\eta+D\) is the other end, the
function \(u(s)=\ell(x,\eta+sD)\) is twice continuously differentiable
on an open interval containing \([0,1]\), with
\(u'(s)=\sum D_\alpha\partial_\alpha\ell(x,\eta+sD)\) and
\(u''(s)=\sum D_\alpha D_\beta\partial_\alpha\partial_\beta\ell(x,\eta+sD)\).
Taylor's formula with integral remainder,
\(u(1)=u(0)+u'(0)+\int_0^1(1-s)u''(s)\,ds\), gives
\[
  \ell(x,\eta)-\ell(x,\eta+D)=-\sum_\alpha D_\alpha\partial_\alpha\ell(x,\eta)-R_{\eta}(x),
\]
where
\[
  R_\eta(x)=\int_0^1(1-s)\sum D_\alpha D_\beta\,\partial_\alpha\partial_\beta\ell(x,\eta+sD)\,ds .
\]
By (R3), \(f(x,\eta)|R_\eta(x)|\leq\tfrac12\|\Delta\|_1^2K(x)\), and by
(R2)
\[
  f(x,\eta)\,|D_\alpha\partial_\alpha\ell(x,\eta)|=|\Delta_\alpha|\,|\partial_\alpha f(x,\eta)|\leq|\Delta_\alpha|\,N(x).
\] So \(f(\cdot,\eta)\bigl(\ell(\cdot,\eta)-\ell(\cdot,\eta+D)\bigr)\)
is \(\lambda\)-integrable. Its integral is the information of the pair
\((\mu_\eta,\mu_{\eta+D})\), which is therefore finite, and by Step~1 at
\(\eta\in B^\circ\) it equals
\[
  -\int f(x,\eta)R_\eta(x)\,d\lambda(x)
  =-\sum\Delta_\alpha\Delta_\beta\,c^{\eta}_{\alpha\beta}(\Delta),
\]
where
\[
  c^{\eta}_{\alpha\beta}(\Delta)=\int_{[0,1]\times X}(1-s)f(x,\eta)\,\partial_\alpha\partial_\beta\ell(x,\eta+sD)\,d(s,x),
\]
using \(D_\alpha D_\beta=\Delta_\alpha\Delta_\beta\). The last step is
Fubini's theorem (\cref{thm:p15-fubini}) for Lebesgue measure on
\([0,1]\) and \(\lambda\), both \(\sigma\)-finite. The integrand is
product measurable: it is continuous in \(s\) for fixed \(x\) and
measurable in \(x\) for fixed \(s\), so it is the pointwise limit of
the product-measurable functions obtained by replacing \(s\) with
\(\lceil ns\rceil/n\). It is bounded by \(K(x)\), whose integral over
\([0,1]\times X\) is \(\int K\,d\lambda<\infty\).

\emph{Step 4: the limits.} For \(\eta=\theta\), \(D=\Delta\), the
information is \(I(\theta;\theta+\Delta)\). For \(\eta=\theta+\Delta\),
\(D=-\Delta\), it is \(I(\theta+\Delta;\theta)\), and
\(\eta+sD=\theta+(1-s)\Delta\). Let \(\Delta^{(n)}\to0\) with
\(0<|\Delta^{(n)}|<r\). In both cases the integrand of
\(c^{\eta}_{\alpha\beta}(\Delta^{(n)})\) converges pointwise on
\([0,1]\times X\) to \((1-s)f(x,\theta)\partial_\alpha\partial_\beta\ell(x,\theta)\),
by continuity of \(f\) and of the second derivatives of \(\ell\) in the
parameter. It is bounded by \(K(x)\) by (R3), since both parameter
points lie in \(B\). Dominated convergence gives
\[
  c^{\eta}_{\alpha\beta}(\Delta^{(n)})\to\int_0^1(1-s)\,ds\int f(x,\theta)\partial_\alpha\partial_\beta\ell(x,\theta)\,d\lambda(x)=-\tfrac12g_{\alpha\beta},
\]
by Step~2. Put
\(\epsilon^\eta_{\alpha\beta}(\Delta)=-c^\eta_{\alpha\beta}(\Delta)-\tfrac12g_{\alpha\beta}\).
Then the information equals
\(\tfrac12\sum g_{\alpha\beta}\Delta_\alpha\Delta_\beta+\sum\epsilon^\eta_{\alpha\beta}(\Delta)\Delta_\alpha\Delta_\beta\),
and
\[
  \Bigl|\sum\epsilon^\eta_{\alpha\beta}(\Delta)\Delta_\alpha\Delta_\beta\Bigr|
  \leq\|\Delta\|_1^2\max_{\alpha,\beta}|\epsilon^\eta_{\alpha\beta}(\Delta)|
  \leq k|\Delta|^2\max_{\alpha,\beta}|\epsilon^\eta_{\alpha\beta}(\Delta)|.
\]
Along \(\Delta^{(n)}\) the maximum tends to \(0\). The sequence was
arbitrary, so both \(I(\theta;\theta+\Delta)\) and
\(I(\theta+\Delta;\theta)\) differ from
\(\tfrac12\sum g_{\alpha\beta}\Delta_\alpha\Delta_\beta\) by a quantity
whose ratio to \(|\Delta|^2\) tends to \(0\). The first is the
statement for \(I\). By \cref{thm:p15-w02b-p15-045-well-defined}(3),
\(J(\theta,\theta+\Delta)=I(\theta;\theta+\Delta)+I(\theta+\Delta;\theta)\),
which is finite and differs from \(\sum g_{\alpha\beta}\Delta_\alpha\Delta_\beta\)
by the sum of the two remainders. This is the statement for \(J\).
\end{proof}
\end{annexitem}

\begin{annexitem}{annex:p15-w02b-p15-045-normal}{thm:p15-w02b-p15-045-normal}{Equation (2.10): normal populations with a common covariance matrix}
\begin{proof}
Write \(\phi_{a}\) for the density of \(\operatorname{normal}(a,1)\),
\(N_a\) for that law, and \(Z=(Z_1,\dots,Z_k)\) for the coordinate map
on \((\mathbb R^k,\rho)\). By \cref{thm:p15-finite-product-measures},
\(\rho=N_0\otimes\dots\otimes N_0\) is the unique measure on the Borel
sets of \(\mathbb R^k\) with
\(\rho(B_1\times\dots\times B_k)=\prod_jN_0(B_j)\).

(1) By Tonelli and Fubini (\cref{thm:p15-tonelli,thm:p15-fubini}) and
\cref{thm:p15-normal-normalization}, \(\int|z_j|\,d\rho<\infty\),
\(\int z_j\,d\rho=0\), \(\int z_j^2\,d\rho=1\), and for \(j\neq l\)
\(\int z_jz_l\,d\rho=\int z_j\,dN_0\int z_l\,dN_0=0\). The
\(\alpha\)-th coordinate of \(m_i+Az\) is
\(m_{i,\alpha}+\sum_jA_{\alpha j}z_j\), so its \(\rho\)-integral is
\(m_{i,\alpha}\), and the integral of the product of the centred
\(\alpha\)-th and \(\beta\)-th coordinates is
\(\sum_{j,l}A_{\alpha j}A_{\beta l}\int z_jz_l\,d\rho=\sum_jA_{\alpha j}A_{\beta j}=\Sigma_{\alpha\beta}\).
The transfer argument in part~(2) of the proof of
\cref{thm:p15-w02b-p15-045-induced} (indicators, simple functions,
monotone convergence), applied to the map \(z\mapsto m_i+Az\) and
\(\rho\), and then to positive and negative parts, gives
\(\int h\,d\mu_i=\int h(m_i+Az)\,d\rho(z)\) for every Borel \(h\) for
which the right side is absolutely integrable. With \(h\) the
coordinates and the products of centred coordinates, these are the mean
vector and covariance matrix of \(\mu_i\).

\emph{The transformation.} Let \(T(x)=A^{-1}(x-m_2)\). It is a
bijection of \(\mathbb R^k\) with inverse \(y\mapsto m_2+Ay\), and both
maps are continuous, hence Borel measurable. For a Borel set \(E\),
\(T(E)\) is the inverse image of \(E\) under \(y\mapsto m_2+Ay\), so it
is Borel, and \(E=T^{-1}(T(E))\). Hence \(T^{-1}\) of the Borel
\(\sigma\)-algebra is the whole Borel \(\sigma\)-algebra: \(T\) is
non-singular in the sense of \cref{thm:p15-w02b-p15-045-nonsingular}.
Let \(\nu_i(G)=\mu_i(T^{-1}G)\). Since \(T(m_1+Az)=d+z\) with
\(d=A^{-1}\delta\), and \(T(m_2+Az)=z\), \(\nu_1\) is the image of
\(\rho\) under \(z\mapsto d+z\) and \(\nu_2=\rho\). On a box,
\(\nu_1(B_1\times\dots\times B_k)=\rho\bigl(\prod_j(B_j-d_j)\bigr)=\prod_jN_0(B_j-d_j)=\prod_jN_{d_j}(B_j)\),
the last step by \cref{thm:p15-normal-sums} with \(c=1\) and shift
\(d_j\). By the uniqueness in \cref{thm:p15-finite-product-measures},
\(\nu_1=N_{d_1}\otimes\dots\otimes N_{d_k}\).

(2) For each \(j\), the pair \((N_{d_j},N_0)\) on \(\mathbb R\)
satisfies the standing hypotheses with \(\lambda_j=N_0\), densities
\(f_1^{(j)}=\phi_{d_j}/\phi_0=e^{d_jx-d_j^2/2}\) and \(f_2^{(j)}=1\),
which are positive and finite; that \(f_1^{(j)}\) is a density follows
from the change-of-measure rule (\cref{thm:p15-density-calculus}),
since \(\phi_0\) is the Lebesgue density of \(N_0\). Applying
\cref{thm:p15-w02b-p15-045-additivity} \(k-1\) times, with
\cref{thm:p15-finite-product-measures} to identify the \(k\)-fold
product with a two-fold product of a \((k-1)\)-fold product and one
factor, the pair \((\nu_1,\nu_2)\) satisfies the standing hypotheses.
In particular \(\nu_1\equiv\nu_2\). For a Borel set \(E\),
\(\mu_i(E)=\mu_i(T^{-1}(T(E)))=\nu_i(T(E))\), so \(\mu_1(E)=0\) if and
only if \(\mu_2(E)=0\): \(\mu_1\equiv\mu_2\). Then
\(\lambda=\tfrac12(\mu_1+\mu_2)\) is a probability measure with the same
null sets, and densities with values in \((0,\infty)\) exist, as noted
after \cref{def:p15-w02b-p15-045-pair}.

(3) By \cref{thm:p15-w02b-p15-045-not-metric}(2), with \(a=d_j\),
\(b=0\), and \cref{thm:p15-w02b-p15-045-well-defined}(2) for the change
of dominating measure, each factor pair has \(I(1:2)=I(2:1)=\tfrac12d_j^2\).
Additivity gives \(I'(1:2)=I'(2:1)=\tfrac12\sum_jd_j^2=\tfrac12|d|^2\)
for the pair \((\nu_1,\nu_2)\); by
\cref{thm:p15-w02b-p15-045-well-defined}(2) these values do not depend
on the dominating measure, so they are the reduced quantities of
\cref{def:p15-w02b-p15-045-statistic} for \(T\). By
\cref{thm:p15-w02b-p15-045-nonsingular}(1), \(I(1:2)=I'(1:2)\). The
standing hypotheses, the induced pair and non-singularity are symmetric
in the indices, so the same corollary with the indices exchanged gives
\(I(2:1)=I'(2:1)\). Finally
\[
  |d|^2=\delta^{\mathsf T}(A^{-1})^{\mathsf T}A^{-1}\delta
       =\delta^{\mathsf T}(AA^{\mathsf T})^{-1}\delta
       =\delta^{\mathsf T}\Sigma^{-1}\delta
       =\sum\delta_\alpha\delta_\beta\sigma^{\alpha\beta},
\]
and \(J(1,2)=I(1:2)+I(2:1)\) by
\cref{thm:p15-w02b-p15-045-well-defined}(3).

(4) \(J(1,2)=|d|^2\) and \(d=A^{-1}(m_1-m_2)=A^{-1}m_1-A^{-1}m_2\).
\end{proof}
\end{annexitem}

\begin{annexitem}{annex:p15-w02b-p15-045-bayes}{thm:p15-w02b-p15-045-bayes}{The Bayes reading of the information in an observation}
\begin{proof}
The joint model is the probability measure \(\Pi\) on
\(\{1,2\}\times X\) with \(\Pi(\{i\}\times E)=\alpha_i\mu_i(E)\). The
marginal law of \(x\) is \(m=\alpha_1\mu_1+\alpha_2\mu_2\), with
\(\lambda\)-density \(\alpha_1f_1+\alpha_2f_2\), which is finite and
positive \([\lambda]\). A version of the posterior probability of \(H_i\)
given \(x\) is an \(\mathbf S\)-measurable \(q_i\) with
\(\Pi(\{i\}\times E)=\int_Eq_i\,dm\) for every \(E\in\mathbf S\). For
\(q_i=\alpha_if_i/(\alpha_1f_1+\alpha_2f_2)\), the change-of-measure
rule gives
\[
  \int_Eq_i\,dm=\int_E\frac{\alpha_if_i}{\alpha_1f_1+\alpha_2f_2}\,(\alpha_1f_1+\alpha_2f_2)\,d\lambda
  =\alpha_i\mu_i(E).
\]
Off a \(\lambda\)-null set both \(\alpha_if_i\) are finite and positive,
so \(q_i\in(0,\infty)\) and \(q_1+q_2=1\). Hence \(q_i\in(0,1)\), and
\(q_1/q_2=\alpha_1f_1/(\alpha_2f_2)\). Taking logarithms gives the
displayed identity \([\lambda]\). Integrating it against \(\mu_1\) gives
\(I(1:2)=\int[\log(q_1/q_2)-\log(\alpha_1/\alpha_2)]\,d\mu_1\), using
\cref{thm:p15-w02b-p15-045-well-defined}(2). This is the mean, under
\(H_1\), of posterior log odds minus prior log odds. The left side is
defined without reference to \(\alpha_1,\alpha_2\). The last assertion
is \cref{thm:p15-w02b-p15-045-footnote}(1) and (3).
\end{proof}
\end{annexitem}

\begin{annexitem}{annex:p15-w02b-p15-045-not-metric}{thm:p15-w02b-p15-045-not-metric}{Neither \(I(1:2)\) nor \(J(1,2)\) is a metric}
\begin{proof}
In both parts the dominating measure is \(\sigma\)-finite with the same
null sets as the pair, which is allowed by
\cref{thm:p15-w02b-p15-045-well-defined}(2).

(1) Counting measure on \(\{0,1\}\) gives densities
\(f_1=(\tfrac12,\tfrac12)\) and \(f_2=(\tfrac14,\tfrac34)\), all
positive. Then
\[
  I(1:2)=\tfrac12\log2+\tfrac12\log\tfrac23=\tfrac12\log\tfrac43,
\]
\[
  I(2:1)=\tfrac14\log\tfrac12+\tfrac34\log\tfrac32
        =\tfrac14\bigl(3\log3-4\log2\bigr)=\tfrac14\log\tfrac{27}{16}.
\]
Numerically \(\tfrac12\log\tfrac43=0.14384\ldots\) and
\(\tfrac14\log\tfrac{27}{16}=0.13081\ldots\), which differ.

(2) The normal densities \(f_a(x)=(2\pi)^{-1/2}e^{-(x-a)^2/2}\) are
positive and finite everywhere, so \(\mu_a\equiv\mu_b\). Also
\[
  \log\frac{f_a(x)}{f_b(x)}=\frac{(x-b)^2-(x-a)^2}2=(a-b)x-\frac{a^2-b^2}2 .
\]
This is \(\mu_a\)-integrable, since \(\int|x|\,d\mu_a<\infty\), and
\(\int x\,d\mu_a=a\) (\cref{thm:p15-normal-normalization}). Hence
\[
  I(1:2)=(a-b)a-\frac{a^2-b^2}2=\frac{(a-b)^2}2 .
\]
Exchanging \(a\) and \(b\) gives the same value for \(I(2:1)\), so
\(J(1,2)=(a-b)^2\) by \cref{thm:p15-w02b-p15-045-well-defined}(3). With
\(a,b,c=0,1,2\) the three divergences are \(J(\mu_0,\mu_2)=4\) and
\(J(\mu_0,\mu_1)=J(\mu_1,\mu_2)=1\), and \(4>1+1\). For \(I\) the
values are \(2\), \(\tfrac12\) and \(\tfrac12\), and \(2>\tfrac12+\tfrac12\).
\end{proof}
\end{annexitem}

\begin{annexitem}{annex:p15-w02b-p15-045-not-utility}{thm:p15-w02b-p15-045-not-utility}{More information can have lower expected utility}
\begin{proof}
(1) The image of \(\mu_i=\mu_i^{(1)}\otimes\mu_i^{(2)}\) under the first
coordinate is \(\mu_i^{(1)}\), since
\(\mu_i(T^{-1}\{t\})=\mu_i(\{t\}\times X_2)=\mu_i^{(1)}(\{t\})\).
Likewise the image under the second coordinate is \(\mu_i^{(2)}\). With
counting measure, which is allowed by
\cref{thm:p15-w02b-p15-045-well-defined}(2), the reduced informations
are
\[
  I_x(1:2)=\tfrac12\log\frac{1/2}{999/1000}+\tfrac12\log\frac{1/2}{1/1000}
          =\tfrac12\log\frac{500}{999}+\tfrac12\log500,
\]
which equals \(\tfrac12\log(250000/999)\), and
\[
  I_y(1:2)=\tfrac9{10}\log9+\tfrac1{10}\log\tfrac19=\tfrac45\log9 .
\]
Numerically \(\tfrac12\log(250000/999)=2.7612\ldots\) and
\(\tfrac45\log9=1.7577\ldots\).

(2) A nonrandomized rule based on a statistic \(S\) with values in
\(\{0,1\}\) is a map \(\delta\) from \(\{0,1\}\) to the two acts. Let
\(\nu_i\) be the law of \(S\) under \(H_i\). The expected utility of
\(\delta\) is
\[
  U(\delta)=\tfrac12\sum_{t\in\{0,1\}}\Bigl(\nu_1(\{t\})\,[\delta(t)=a_1]
                       +\nu_2(\{t\})\,[\delta(t)=a_2]\Bigr),
\]
where \(a_i\) is the act ``assert \(H_i\)''. Hence
\[
  U(\delta)\leq\tfrac12\sum_t\max\bigl(\nu_1(\{t\}),\nu_2(\{t\})\bigr),
\]
where \([\cdot]\) is \(1\) for a true condition and \(0\) otherwise.
Equality holds for the rule that asserts the hypothesis with the larger
\(\nu_i(\{t\})\). A randomized rule's expected utility is an average of
these values, so it obeys the same bound. For \(S=T\), the bound is
\(\tfrac12(\tfrac{999}{1000}+\tfrac12)=\tfrac{1499}{2000}\). For \(S=T'\)
it is \(\tfrac12(\tfrac9{10}+\tfrac9{10})=\tfrac9{10}\).
\end{proof}
\end{annexitem}

\begin{annexitem}{annex:p15-w02b-p15-045-footnote}{thm:p15-w02b-p15-045-footnote}{Positive densities, infinite information}
\begin{proof}
For \(0<t<1\), \(0<e^{-1/t}<e^{-1}\), so \(0<k^{-1}<e^{-1}\) and
\(k\in(e,\infty)\). \(f_2\) is continuous, positive and finite on
\((0,1)\), and \(\int_0^1f_2\,d\lambda=k\cdot k^{-1}=1\). So \(\mu_2\) is
a probability measure with a positive density with respect to
\(\lambda=\mu_1\). Hence \(\mu_1\equiv\mu_2\), \(\lambda\) is a
probability measure with \(\lambda\equiv\{\mu_1,\mu_2\}\), and the
standing hypotheses hold with \(f_1=1\).

(1) \(\log(f_1/f_2)(x)=1/x-\log k\). Its negative part is at most
\(\log k\), and
\(\int_0^1(1/x-\log k)^+\,dx\geq\int_0^{1/\log k}(1/x-\log k)\,dx=\infty\),
since \(\int_0^\delta dx/x=\infty\) for every \(\delta>0\). So
\(I(1:2)=\infty\).

(2) \(\log(f_2/f_1)(x)=\log k-1/x\), so
\(I(2:1)=\int_0^1ke^{-1/x}(\log k-1/x)\,dx\). With \(u=1/x\geq1\),
\(x^{-1}e^{-1/x}=ue^{-u}\leq e^{-1}\). So
\(0\leq\int_0^1x^{-1}e^{-1/x}\,dx\leq e^{-1}\) is finite, and (E1) gives
\(I(2:1)=\log k-k\int_0^1x^{-1}e^{-1/x}\,dx\leq\log k<\infty\).
\(I(2:1)\geq0\) by \cref{thm:p15-w02b-p15-045-nonnegative}. The decimal
values in the statement come from numerical quadrature of
\(\int_0^1e^{-1/t}\,dt\) and \(\int_0^1t^{-1}e^{-1/t}\,dt\). They are
not used in any argument.

(3) By \cref{thm:p15-w02b-p15-045-bayes}, a version of the posterior is
\(P(H_1\mid x)=\alpha_1/(\alpha_1+\alpha_2ke^{-1/x})\) and
\(P(H_2\mid x)=1-P(H_1\mid x)\). Both lie in \((0,1)\) for every
\(x\in(0,1)\), because every term is positive and finite.
\end{proof}
\end{annexitem}

\begin{annexitem}{annex:p15-w02b-p15-045-support}{thm:p15-w02b-p15-045-support}{Support mismatch: the extended values as limits}
\begin{proof}
(1) Let \(\bar\mu=\tfrac12(\mu_1+\mu_2)\), a probability measure. For
\(A\in\mathbf S\), \(\bar\mu(A)=0\) if and only if
\(\mu_1(A)=\mu_2(A)=0\). Since both coefficients \(1-\varepsilon\) and
\(\varepsilon\) are positive, this happens if and only if
\(\bar\mu_1^\varepsilon(A)=0\), and if and only if
\(\bar\mu_2^\varepsilon(A)=0\). So
\(\bar\mu_1^\varepsilon\equiv\bar\mu_2^\varepsilon\), and
\(\lambda=\bar\mu\) is a probability measure with
\(\lambda\equiv\{\bar\mu_1^\varepsilon,\bar\mu_2^\varepsilon\}\).
Densities with values in \((0,\infty)\) exist, as noted after
\cref{def:p15-w02b-p15-045-pair}, so the standing hypotheses hold. Also
\(|\bar\mu_1^\varepsilon(A)-\mu_1(A)|=\varepsilon|\mu_2(A)-\mu_1(A)|\leq\varepsilon\),
and likewise for \(\bar\mu_2^\varepsilon\).

Let \(F=X\setminus E\). By \cref{thm:p15-w02b-p15-045-well-defined}(4),
\[
  \bar I^\varepsilon(1:2)=\bar\mu_1^\varepsilon(E)\,\bar I^\varepsilon_{1:2}(E)
                       +\bar\mu_1^\varepsilon(F)\,\bar I^\varepsilon_{1:2}(F).
\]
Since \(\mu_2(E)=0\), \(\bar\mu_1^\varepsilon(E)=(1-\varepsilon)\mu_1(E)>0\)
and \(\bar\mu_2^\varepsilon(E)=\varepsilon\mu_1(E)>0\); hence
\(\lambda(E)>0\), and \cref{thm:p15-w02b-p15-045-event-bound} gives
\(\bar I^\varepsilon_{1:2}(E)\geq\log((1-\varepsilon)/\varepsilon)\).
If \(t=\bar\mu_1^\varepsilon(F)>0\), then \(\lambda(F)>0\), and the same
lemma with \(\bar\mu_2^\varepsilon(F)\leq1\) gives
\(\bar I^\varepsilon_{1:2}(F)\geq\log t\). The second term is then at
least \(t\log t\geq-1/e\). If \(t=0\), the second term is \(0\). Adding
gives the bound, a sum in \((-\infty,\infty]\). As \(\varepsilon\to0\),
\((1-\varepsilon)\log((1-\varepsilon)/\varepsilon)\to\infty\), and
\(\mu_1(E)>0\) does not depend on \(\varepsilon\).

(2) \emph{The first direction.} If \(\mu_1(A)=0\), then \(\mu_2(A)=0\) because \(\mu_2\ll\mu_1\), so
\(\mu_2^\varepsilon(A)=0\). If \(\mu_2^\varepsilon(A)=0\), then
\(\varepsilon\mu_1(A)=0\). So \(\mu_1\equiv\mu_2^\varepsilon\), and
\(\lambda=\mu_1\) is a probability measure with the same null sets. The
standing hypotheses hold, since densities with values in \((0,\infty)\)
exist, as noted after \cref{def:p15-w02b-p15-045-pair}. Let
\(F=X\setminus E\). By \cref{thm:p15-w02b-p15-045-well-defined}(4),
\(I^\varepsilon(1:2)=\mu_1(E)I^\varepsilon_{1:2}(E)+\mu_1(F)I^\varepsilon_{1:2}(F)\).
Since \(\mu_2(E)=0\), \(\mu_2^\varepsilon(E)=\varepsilon\mu_1(E)>0\), and
\cref{thm:p15-w02b-p15-045-event-bound} gives
\(I^\varepsilon_{1:2}(E)\geq\log(\mu_1(E)/(\varepsilon\mu_1(E)))=\log(1/\varepsilon)\).
If \(\mu_1(F)>0\), the same lemma and \(\mu_2^\varepsilon(F)\leq1\) give
\(I^\varepsilon_{1:2}(F)\geq\log(\mu_1(F)/\mu_2^\varepsilon(F))\geq\log\mu_1(F)\).
If \(\mu_1(F)=0\), that term is \(0\). Adding gives the bound. As
\(\varepsilon\to0\) the first term tends to \(\infty\) and the second
does not depend on \(\varepsilon\).

\emph{The second direction.} Let \(g\) be a version of
\(d\mu_2/d\mu_1\) with \(0\leq g<\infty\). Then
\(g_\varepsilon=(1-\varepsilon)g+\varepsilon\) is a version of
\(d\mu_2^\varepsilon/d\mu_1\) with values in \([\varepsilon,\infty)\),
and with \(\lambda=\mu_1\), \(f_1=1\), \(f_2=g_\varepsilon\),
\[
  I^\varepsilon(2:1)=\int g_\varepsilon\log g_\varepsilon\,d\mu_1=\int\varphi(g_\varepsilon)\,d\mu_1,
\]
where \(\varphi\) is as in \cref{thm:p15-w02b-p15-045-extension}. By
\cref{thm:p15-w02b-p15-045-extension}(1) with the indices exchanged,
\(I(2:1)=\int\varphi(g)\,d\mu_1\in[0,\infty]\), and \(I(2:1)>0\)
because \(\mu_2\neq\mu_1\).

\(\varphi\) is convex on \([0,\infty)\). On \((0,\infty)\) its
derivative \(1+\log t\) is increasing, so by the mean value theorem the
slope of the chord of \(\varphi\) over \([s,t]\) is at most the slope
over \([t,u]\) for \(0<s<t<u\), which is convexity on \((0,\infty)\);
continuity of \(\varphi\) at \(0\) extends the chord inequality to
\([0,\infty)\). Hence, pointwise,
\[
  \varphi(g_\varepsilon)=\varphi\bigl((1-\varepsilon)g+\varepsilon\cdot1\bigr)
  \leq(1-\varepsilon)\varphi(g)+\varepsilon\varphi(1)=(1-\varepsilon)\varphi(g).
\]
Write \(u=\varphi(g_\varepsilon)\) and \(v=(1-\varepsilon)\varphi(g)\).
Both have negative parts at most \(1/e\), hence \(\mu_1\)-integrable.
From \(u\leq v\), \(u^++v^-\leq v^++u^-\) pointwise; integrating
nonnegative functions and subtracting the finite quantities
\(\int u^-\,d\mu_1\) and \(\int v^-\,d\mu_1\) gives
\(\int u\,d\mu_1\leq\int v\,d\mu_1\), that is,
\(I^\varepsilon(2:1)\leq(1-\varepsilon)I(2:1)\).

For the limit, let \(\varepsilon_n\to0\). Pointwise
\(\varphi(g_{\varepsilon_n})+1/e\to\varphi(g)+1/e\) by continuity of
\(\varphi\), and both are nonnegative. Fatou's inequality
(\cref{thm:p15-nonnegative-integral-laws}) and \(\mu_1(X)=1\) give
\(\liminf_n I^{\varepsilon_n}(2:1)\geq I(2:1)\). If \(I(2:1)=\infty\)
this gives \(I^{\varepsilon_n}(2:1)\to\infty\). If \(I(2:1)<\infty\), the
upper bound gives
\(\limsup_nI^{\varepsilon_n}(2:1)\leq\lim_n(1-\varepsilon_n)I(2:1)=I(2:1)\).
In both cases \(I^{\varepsilon_n}(2:1)\to I(2:1)\), and the sequence was
arbitrary.

(3) Here \(\mu_2\ll\mu_1\), with \(d\mu_2/d\mu_1=2\cdot1_{(0,1)}\), and
\(E=[1,2)\) has \(\mu_1(E)=\tfrac12>0=\mu_2(E)\). With \(\lambda=\mu_1\),
\(f_1=1\) and \(f_2^\varepsilon=d\mu_2^\varepsilon/d\mu_1\), which is
\((1-\varepsilon)2+\varepsilon=2-\varepsilon\) on \((0,1)\) and
\(\varepsilon\) on \([1,2)\). Each piece has \(\mu_1\)-mass \(\tfrac12\).
Hence
\[
  I^\varepsilon(1:2)=\tfrac12\log\frac1{2-\varepsilon}+\tfrac12\log\frac1\varepsilon
                   =-\tfrac12\log\bigl(\varepsilon(2-\varepsilon)\bigr),
\]
\[
  I^\varepsilon(2:1)=\int f_2^\varepsilon\log f_2^\varepsilon\,d\mu_1
   =\tfrac12(2-\varepsilon)\log(2-\varepsilon)+\tfrac12\varepsilon\log\varepsilon .
\]
As \(\varepsilon\to0\), \(-\tfrac12\log(\varepsilon(2-\varepsilon))\to\infty\).
Since \(\varepsilon\log\varepsilon\to0\), the second quantity tends to
\(\log2\). For the unmixed pair,
\(I(2:1)=\int\log(d\mu_2/d\mu_1)\,d\mu_2=\log2\) by
\cref{def:p15-w02b-p15-045-extension}(1), because
\(d\mu_2/d\mu_1=2\) \([\mu_2]\). This agrees with the limit asserted
in (2).
\end{proof}
\end{annexitem}

\begin{annexitem}{annex:p15-w02b-p15-045-loss}{thm:p15-w02b-p15-045-loss}{Information loss by a statistic that is not sufficient}
\begin{proof}
By \cref{thm:p15-w02b-p15-045-additivity},
\(I(1:2)=I_x(1:2)+I_y(1:2)\). The induced pair of \(T\) is
\((\mu_1^{(1)},\mu_2^{(1)})\), as in the proof of
\cref{thm:p15-w02b-p15-045-not-utility}, so \(I'(1:2)=I_x(1:2)\). Since
\(I_y(1:2)=\tfrac45\log9\) is positive and finite, \(I'(1:2)<I(1:2)\),
and both are finite. By \cref{thm:p15-w02b-p15-045-reduction}(3), \(T\)
is not sufficient. Directly, \(f_1/f_2\) at \((x,0)\) and \((x,1)\) is
\(r(x)\cdot9\) and \(r(x)/9\) with \(r(x)>0\). These differ, so the ratio
is not a function of \(x\), and every point carries positive mass. The
numerical value of the sum is \(2.7612\ldots+1.7577\ldots=4.5190\ldots\).
\end{proof}
\end{annexitem}

\begin{annexitem}{annex:p15-w02b-p15-045-infinite-equality}{thm:p15-w02b-p15-045-infinite-equality}{Equal infinite information without sufficiency}
\begin{proof}
Use \(\lambda_1\) Lebesgue measure on \((0,1)\) and \(\lambda_2\) the
uniform probability on \(\{0,1\}\). Both factor pairs satisfy the
standing hypotheses (\cref{thm:p15-w02b-p15-045-footnote}, and positive
masses on \(\{0,1\}\)). By \cref{thm:p15-w02b-p15-045-additivity} the
product pair does too, with
\(\log(f_1/f_2)(x,y)=a(x)+b(y)\), where \(a(x)=1/x-\log k\),
\(b(0)=\log9\) and \(b(1)=-\log9\).

(1) By additivity, \(I(1:2)=I_x(1:2)+I_y(1:2)=\infty+\tfrac45\log9=\infty\).
The induced pair of \(T\) is \((\mu_1^{(1)},\mu_2^{(1)})\) with
\(\gamma=\lambda_1\) and \(g_i=f_i^{(1)}\). So
\(I'(1:2)=I_x(1:2)=\infty\). Since \(J=I(1:2)+I(2:1)\) and
\(J'=I'(1:2)+I'(2:1)\) (\cref{thm:p15-w02b-p15-045-well-defined}(3)),
both are \(\infty\).

Suppose \(T\) were sufficient. By
\cref{thm:p15-w02b-p15-045-density-criterion} there would be a
\(T^{-1}(\mathbf T)\)-measurable \(W\) with \(f_1/f_2=W\) \([\lambda]\).
Every member of \(T^{-1}(\mathbf T)\) has the form \(A\times\{0,1\}\).
Taking the event \(\{W=W(x,0)\}\) shows \(W(x,0)=W(x,1)\) for every
\(x\). But \((f_1/f_2)(x,0)=9e^{a(x)}\neq e^{a(x)}/9=(f_1/f_2)(x,1)\).
So for every \(x\), at least one of \((x,0),(x,1)\) lies in the
measurable set \(D=\{f_1/f_2\neq W\}\). Tonelli gives
\(\lambda(D)=\int\lambda_2(D_x)\,d\lambda_1(x)\geq\tfrac12\), which
contradicts \(\lambda(D)=0\). So \(T\) is not sufficient.

(2) Let \(G=(\tfrac12,1)\), so \(T^{-1}G=G\times\{0,1\}\) and
\(\mu_1(T^{-1}G)=\nu_1(G)=\tfrac12\). On \(G\), \(a\) is bounded, so
\(I'_{1:2}(G)=2\int_{1/2}^1(1/x-\log k)\,dx=2\log2-\log k\) is finite.
As in the proof of \cref{thm:p15-w02b-p15-045-additivity}, (E1) and
Tonelli split the integral of \(f_1(a+b)\) over \(G\times\{0,1\}\) into
\(\int_Ga\,dx\cdot1+\mu_1^{(1)}(G)\,I_y(1:2)\). Hence
\[
  I_{1:2}(T^{-1}G)=2\Bigl(\int_{1/2}^1a\,dx+\tfrac12\cdot\tfrac45\log9\Bigr)
                 =I'_{1:2}(G)+\tfrac45\log9\neq I'_{1:2}(G).
\]

(3) Also \(\mu_2(T^{-1}G)=\nu_2(G)=\int_{1/2}^1ke^{-1/x}\,dx>0\), and
\(\log(f_2/f_1)(x,y)=-a(x)-b(y)\). On \(G\) the function
\(f_2^{(1)}\cdot(-a)\) is bounded, so
\[
  I'_{2:1}(G)=\frac1{\nu_2(G)}\int_{1/2}^1ke^{-1/x}\bigl(\log k-1/x\bigr)\,dx
\]
is finite. The integrand \(f_2\cdot(-a-b)\) is bounded on
\(G\times\{0,1\}\), so the integral over \(G\times\{0,1\}\) splits by
Tonelli into \(\int_Gf_2^{(1)}(-a)\,dx\cdot1+\nu_2(G)\,I_y(2:1)\), where
\[
  I_y(2:1)=\tfrac1{10}\log\tfrac19+\tfrac9{10}\log9=\tfrac45\log9 .
\]
Dividing by \(\mu_2(T^{-1}G)=\nu_2(G)\) gives
\(I_{2:1}(T^{-1}G)=I'_{2:1}(G)+\tfrac45\log9\). Adding this to (2),
\[
  J_{12}(T^{-1}G)=I'_{1:2}(G)+I'_{2:1}(G)+\tfrac85\log9
                 =J'_{12}(G)+\tfrac85\log9 ,
\]
with \(J'_{12}(G)\) finite, so \(J_{12}(T^{-1}G)\neq J'_{12}(G)\).
\end{proof}
\end{annexitem}

\begin{annexitem}{annex:p15-w02b-p15-045-critical}{thm:p15-w02b-p15-045-critical}{Section 6(1): comparing two critical regions}
\begin{proof}
(1) Every point of \(X\) has positive mass under both measures, so a set
is null for \(\mu_1\), for \(\mu_2\) or for
\(\tfrac12(\mu_1+\mu_2)\) exactly when it is empty. Hence
\(\mu_1\equiv\mu_2\), and the standing hypotheses hold with the
probability measure \(\tfrac12(\mu_1+\mu_2)\). By
\cref{thm:p15-w02b-p15-045-well-defined}(2) the quantities may be
computed with counting measure on \(X\), which is \(\sigma\)-finite
because \(X\) is countable and has the same null sets. The densities are
then the point masses, and an integral over a finite set is the sum over
its points. Since \(\mu_1(F)=\alpha>0\),
\[
  I_{1:2}(F)=\frac1\alpha\Bigl(\frac\alpha2\log\frac{\alpha/2}{pc}
            +\frac\alpha2\log\frac{\alpha/2}{(1-p)c}\Bigr),
\]
which is the displayed formula. Since \(\mu_1(G)=\mu_1(\{0\})>0\),
\[
  I_{1:2}(G)=\frac1{\mu_1(\{0\})}\,\mu_1(\{0\})\log\frac{\mu_1(\{0\})}{\mu_2(\{0\})}
           =\log\frac{1-\alpha}{1-c}.
\]

(2) Expanding the logarithms,
\[
  I_{1:2}(F)=\log\frac\alpha{2c}-\frac12\log p-\frac12\log(1-p)
           =\log\frac\alpha{2c}-\frac12\log\bigl(p(1-p)\bigr),
\]
and subtracting \(I_{1:2}(G)=\log((1-\alpha)/(1-c))\) gives the
stated difference.

(3) By (2), \(I_{1:2}(F)>I_{1:2}(G)\) if and only if
\(\alpha(1-c)/(2c(1-\alpha))>\sqrt{p(1-p)}\), since the logarithm is
strictly increasing. Multiplying by \(2c(1-\alpha)>0\) gives the stated
condition. Next, \(4p(1-p)=1-(2p-1)^2\), so
\(2\sqrt{p(1-p)}\leq1\), with equality if and only if \(p=\tfrac12\).
Also \(\alpha(1-c)-c(1-\alpha)=\alpha-c\).
If \(c<\alpha\), then
\(\alpha(1-c)>c(1-\alpha)\geq2c(1-\alpha)\sqrt{p(1-p)}\).
If \(c=\alpha\) and \(p\neq\tfrac12\), then
\(\alpha(1-c)=c(1-\alpha)>2c(1-\alpha)\sqrt{p(1-p)}\).
If \(p=\tfrac12\), the condition reads \(\alpha(1-c)>c(1-\alpha)\), that
is, \(c<\alpha\); so it fails when \(c\geq\alpha\).
\end{proof}
\end{annexitem}
